\chapter{Modern \cpp{} for \maya{}}
\label{ch:cpp}

This chapter is a self-contained course in the slice of \cpp{} that
\maya{} uses. It assumes no \cpp{} background. If you have written
\cpp{}11 or \cpp{}14 it will still be useful, because \maya{} leans
on features that arrived in \cpp{}20, \cpp{}23, and \cpp{}26; some of
the patterns will be new even to experienced \cpp{} programmers.

Read it linearly the first time. The sections build on each other:
later ones (templates, concepts, variants) lean on earlier ones
(values, references, classes). When you finish, you will be able to
read every line of \maya{} source code.

\section{What a \cpp{} program is}
\label{sec:cpp:program}

A \cpp{} program is one or more \textbf{source files} compiled
together into an \textbf{executable}. The smallest possible one fits
on one line, although by convention we spread it out:

\begin{lstlisting}
#include <print>

int main() {
    std::println("Hello, world");
}
\end{lstlisting}

Four things are happening, and each deserves a paragraph.

\paragraph{The preprocessor runs first.}
Before the compiler proper looks at your code, a small program called
the \textbf{preprocessor} does textual substitution. Lines beginning
with \code{\#} are its instructions. \code{\#include <print>}
literally means ``find the file named \file{print} in the standard
library's directories and paste its contents here''. After the
preprocessor finishes, your source file has grown to several
thousand lines, all the standard-library declarations now visible.

\paragraph{The compiler reads what is left.}
The expanded text is parsed, type-checked, and translated into
\textbf{object code} --- machine instructions in a not-yet-runnable
format. Each source file becomes one \emph{object file}, also called
a \textbf{translation unit}. A small project has one; \maya{} has
several hundred.

\paragraph{The linker stitches object files together.}
The \textbf{linker} takes all object files plus any libraries you
depend on and produces a single executable. Its main job is
resolving \textbf{symbols}: when one file uses a function defined in
another file, the linker matches the use to the definition.

\paragraph{The operating system runs the result.}
When you type \code{./hello}, the OS loads the executable into
memory and jumps to a specific named function: \code{main}. The
integer \code{main} returns is the process's \emph{exit status}
(zero = success). The whole program ends when \code{main} returns.

\begin{idea}
The preprocessor pastes text. The compiler turns text into object
files. The linker joins object files into an executable. The OS runs
the executable from \code{main}.
\end{idea}

\subsection{Headers, source files, and the one-definition rule}

A \textbf{header file} (suffix \file{.hpp} or \file{.h}) is a file
you \code{\#include} into other files. It typically contains
\emph{declarations} --- statements about what exists, but not how it
works. A \textbf{source file} (suffix \file{.cpp}) typically contains
\emph{definitions} --- the actual function bodies and data.

\begin{lstlisting}
// math.hpp --- declaration
int add(int a, int b);
\end{lstlisting}

\begin{lstlisting}
// math.cpp --- definition
#include "math.hpp"
int add(int a, int b) { return a + b; }
\end{lstlisting}

\begin{lstlisting}
// main.cpp --- user
#include "math.hpp"
int main() { return add(2, 3); }
\end{lstlisting}

Every use of \code{add} in any source file sees the same declaration
(from \file{math.hpp}), and there is exactly one definition in the
whole program (in \file{math.cpp}). That is the \textbf{One Definition
Rule}, ODR: every entity must have exactly one definition across all
translation units linked together.

\maya{} is a \emph{header-mostly} library. Almost all its code lives
in headers, with most functions and classes defined \code{inline} or
as templates. This is unusual and has trade-offs (compile times can
be long), but it lets the compiler aggressively inline and constant-
fold across the library boundary. We discuss the trade-off in
Chapter~\ref{ch:build}.

\subsection{\code{\#pragma once} and include guards}

If two headers include the same third header, the preprocessor would
paste it twice, producing duplicate declarations. The traditional
fix is an \textbf{include guard}:

\begin{lstlisting}
// math.hpp
#ifndef MATH_HPP
#define MATH_HPP
int add(int a, int b);
#endif
\end{lstlisting}

The second inclusion finds \code{MATH\_HPP} already defined and skips
the contents. Modern code uses the shorter and equivalent:

\begin{lstlisting}
// math.hpp
#pragma once
int add(int a, int b);
\end{lstlisting}

Every file in \maya{} starts with \code{\#pragma once}.

\section{Values, references, pointers, \code{const}}
\label{sec:cpp:memory}

\cpp{} forces you to be precise about \emph{where} a piece of data
lives and \emph{who can change it}. This precision is what makes
\cpp{} fast and what makes beginners stumble. We will walk it slowly.

\subsection{Values}

A \textbf{value} is just data: \code{int x = 42;} creates a variable
\code{x} of type \code{int} holding the value $42$. The variable
owns its storage. When \code{x} goes out of scope (the block ends),
the storage is gone.

\begin{lstlisting}
int x = 42;          // x owns 4 bytes holding 42
int y = x;           // y owns its own 4 bytes, copy of x
y = 100;             // x is still 42
\end{lstlisting}

This is the \textbf{value semantics} model: every assignment is a
copy, every variable is its own thing. The same applies to classes:

\begin{lstlisting}
std::string a = "hello";
std::string b = a;   // b is its own string, copy of a\'s characters
b += " world";
// a == "hello", b == "hello world"
\end{lstlisting}

For a \code{std::string} this is a real copy of every character.
That sounds wasteful and sometimes is; we will see \code{std::move}
shortly as the way to opt out of copies when you no longer need the
original.

\subsection{References}

A \textbf{reference} is an alias for an existing value. It is \emph{not}
a separate variable; it is another name for the same storage.

\begin{lstlisting}
int x = 42;
int& r = x;        // r is an alias for x
r = 100;
// x is now 100
\end{lstlisting}

References must be bound when created and cannot be re-seated.
\code{int\& r;} alone is not legal --- a reference always names
something.

References are the standard way to pass large values into a function
without copying:

\begin{lstlisting}
void print_long(const std::string& s) {
    std::println("{}", s);     // s is the caller\'s string
}                                  // no copy, no destruction
\end{lstlisting}

The \code{const} qualifier on a reference promises ``I will not
modify the thing you let me look at''. \code{const T\&} is the
idiomatic way to receive read-only inputs of any non-trivial type.

\subsection{Pointers}

A \textbf{pointer} is a value that holds an address --- the location
of another value in memory. The syntax \code{T*} means ``pointer to a
\code{T}''. You dereference with \code{*}.

\begin{lstlisting}
int x = 42;
int* p = &x;         // p holds the address of x
*p = 100;            // write through the pointer
// x is now 100
\end{lstlisting}

Unlike references, pointers \emph{can} be re-seated, can be null
(\code{nullptr}), and can point at nothing meaningful (a real source
of bugs). \maya{} avoids raw pointers in user-facing APIs.
Internally, it uses them for things that genuinely model ``maybe
points to something'': linked lists, lookup tables, parent links
in a tree.

When you do see a pointer in \maya{}, prefer the modern alternatives:
\code{std::unique\_ptr<T>} (single owner, automatically destroyed),
\code{std::shared\_ptr<T>} (reference-counted shared owner),
\code{std::optional<T>} (a value that may or may not be present, no
heap allocation), \code{std::span<T>} (a non-owning view of a
contiguous range), \code{std::string\_view} (the same for character
data).

\subsection{\code{const} on values and references}

A \code{const} variable cannot be reassigned:

\begin{lstlisting}
const int n = 5;
// n = 10;          // error: cannot assign to const
\end{lstlisting}

A \code{const} reference promises read-only access:

\begin{lstlisting}
void f(const std::string& s) {
    // s.push_back('!');  // error: s is const here
    std::println("{}", s);
}
\end{lstlisting}

A \code{const} method on a class promises not to modify the object:

\begin{lstlisting}
struct Counter {
    int n = 0;
    int get() const { return n; }   // const method, ok to call on const objects
    void inc()       { ++n; }
};

const Counter c;
c.get();           // ok
// c.inc();        // error: cannot call non-const method on const object
\end{lstlisting}

\code{const}-correctness is a deceptively powerful discipline: it
lets the compiler check at compile time that your read-only paths
really are read-only.

\begin{recap}
A value owns its storage. A reference is an alias; a pointer holds
an address. \code{const} marks read-only at every level: variable,
reference, method.
\end{recap}

\section{Functions, overloads, default arguments}
\label{sec:cpp:functions}

A function takes inputs, may return an output, and may have side
effects:

\begin{lstlisting}
int square(int n) { return n * n; }
void greet(std::string_view name) {
    std::println("hi, {}", name);
}
\end{lstlisting}

The return type goes first, then the name, then the parameters in
parentheses. \code{void} means ``returns nothing''.

\subsection{Parameter passing}

A parameter declared by value (\code{int n}) is a fresh copy. A
parameter declared by reference (\code{const T\&}) shares storage
with the caller's value. The rule of thumb:

\begin{itemize}[leftmargin=*]
  \item Trivial small types (\code{int}, \code{double},
    \code{std::string\_view}, \code{std::span<T>}) --- by value.
    Copying is cheap.
  \item Larger types you only read --- by \code{const} reference.
    No copy.
  \item Types you intend to modify in place --- by non-\code{const}
    reference. The caller sees the change.
  \item Types you want to take ownership of --- by value (and let
    the caller \code{std::move} into the call).
\end{itemize}

\subsection{Overloading}

Multiple functions can share a name if their parameter lists differ.
This is \textbf{overloading}:

\begin{lstlisting}
void log(int n);
void log(double d);
void log(std::string_view s);

log(42);          // calls log(int)
log(3.14);        // calls log(double)
log("hello");     // calls log(string_view)
\end{lstlisting}

The compiler picks the best match. If two are equally good, you get
a compile error and you spell out the choice.

\subsection{Default arguments}

A parameter can have a default:

\begin{lstlisting}
void greet(std::string_view name, std::string_view greeting = "hi") {
    std::println("{}, {}", greeting, name);
}

greet("world");            // "hi, world"
greet("world", "hello");   // "hello, world"
\end{lstlisting}

Defaults are part of the declaration; later overloads of the same
name cannot redeclare them.

\section{Classes, constructors, destructors}
\label{sec:cpp:classes}

A \textbf{class} (introduced with \code{class} or \code{struct} ---
the only difference is the default access level) is a user-defined
type that bundles data and behaviour.

\begin{lstlisting}
struct Point {
    double x;
    double y;

    double length() const {
        return std::sqrt(x*x + y*y);
    }
};

Point p{3.0, 4.0};
std::println("{}", p.length());   // 5
\end{lstlisting}

\code{struct} members are public by default; \code{class} members
are private by default. Most \maya{} types are \code{struct}.

\subsection{Constructors}

A \textbf{constructor} is a special function that runs when an object
is created. It has the same name as the class, no return type, and
may take parameters.

\begin{lstlisting}
struct File {
    FILE* fp;

    explicit File(const char* path) {
        fp = std::fopen(path, "r");
    }
};

File log{"app.log"};   // calls the constructor with "app.log"
\end{lstlisting}

\code{explicit} prevents implicit conversion: without it, you could
write \code{File f = "app.log";} and \cpp{} would silently call the
constructor. Mark single-argument constructors \code{explicit} unless
you genuinely want the implicit form.

\subsection{Member initialiser lists}

The preferred way to set members from a constructor is the
\textbf{initialiser list}, which appears between a colon and the
body:

\begin{lstlisting}
struct File {
    FILE* fp;

    explicit File(const char* path) : fp(std::fopen(path, "r")) {}
};
\end{lstlisting}

Members are initialised in the order they are declared in the class,
\emph{not} in the order they appear in the initialiser list. Modern
compilers warn if you write them in the wrong order; pay attention
to the warning.

\subsection{Destructors}

A \textbf{destructor} is the constructor's counterpart: it runs when
the object is about to be destroyed. Its name is the class name
prefixed with \code{\~{}}.

\begin{lstlisting}
struct File {
    FILE* fp;

    explicit File(const char* path) : fp(std::fopen(path, "r")) {}
    ~File() { if (fp) std::fclose(fp); }
};

void load() {
    File log{"app.log"};   // fopen runs here
    // ... use log.fp ...
}                              // ~File runs here, fclose runs automatically
\end{lstlisting}

The destructor runs at the end of the scope where the object was
created, even if the function returns early, even if an exception is
thrown. This is the foundation of the next idea.

\section{RAII --- the central pattern}
\label{sec:cpp:raii}

\textbf{Resource Acquisition Is Initialization}, or \textbf{RAII}, is
the single most important idiom in \cpp{}. It is the principle that
every resource your program holds --- memory, file descriptor,
terminal mode, mutex lock, network connection --- should be tied to
the lifetime of an object. The constructor acquires; the destructor
releases. You cannot forget to release because the language runs the
destructor for you.

The \code{File} class above is RAII: constructing a \code{File} opens
the file; the destructor closes it. There is no \code{close()}
method to forget to call.

\subsection{A \code{TerminalGuard} for raw mode}

Let us write the RAII type that fixes the bug from
Chapter~\ref{ch:terminal}, where forgetting to restore cooked mode
left the user's shell broken.

\begin{lstlisting}
class TerminalGuard {
    termios cooked_;
public:
    TerminalGuard() {
        tcgetattr(STDIN_FILENO, &cooked_);
        termios raw = cooked_;
        raw.c_lflag &= ~(ICANON | ECHO | ISIG | IEXTEN);
        raw.c_iflag &= ~(IXON | ICRNL | INPCK | ISTRIP);
        raw.c_oflag &= ~(OPOST);
        raw.c_cc[VMIN]  = 0;
        raw.c_cc[VTIME] = 0;
        tcsetattr(STDIN_FILENO, TCSAFLUSH, &raw);
    }

    ~TerminalGuard() {
        tcsetattr(STDIN_FILENO, TCSAFLUSH, &cooked_);
    }

    TerminalGuard(const TerminalGuard&)            = delete;
    TerminalGuard& operator=(const TerminalGuard&) = delete;
};

int main() {
    TerminalGuard guard;   // raw mode on
    // ... run a TUI ...
}                          // raw mode off, automatically
\end{lstlisting}

The two \code{= delete} lines forbid copying. Copying would let two
\code{TerminalGuard} objects fight over the terminal state. We will
see \code{= delete} again shortly.

\maya{}'s actual terminal driver is more elaborate (it also handles
the alt-screen, mouse mode, KKP push, signal restoration on crash),
but the shape is exactly this: a class whose lifetime is the
lifetime of the terminal being in TUI mode.

\begin{idea}
RAII: every resource is owned by an object. Acquire in the
constructor, release in the destructor. Forgetting becomes a
compile-time impossibility, not a runtime risk.
\end{idea}

\subsection{The rule of zero, three, five}

If a class manages a resource, you usually need to define five
special member functions: destructor, copy constructor, copy
assignment, move constructor, move assignment. This is the
\textbf{rule of five}.

If you do not manage any resource yourself --- you only hold other
RAII types as members (\code{std::string}, \code{std::vector},
\code{std::unique\_ptr}) --- you should define \emph{none} of them.
The compiler generates correct versions automatically. This is the
\textbf{rule of zero}, and it is the rule you should follow first.

The \textbf{rule of three} (\cpp{}98) is the pre-move-semantics
version: define destructor, copy constructor, copy assignment
together.

Most \maya{} classes follow the rule of zero. The few that hold raw
resources (the terminal guard, signal handlers, file descriptor
wrappers) follow the rule of five.

\section{\code{auto} and type deduction}
\label{sec:cpp:auto}

Writing out the type of every variable is fine for \code{int} and
\code{double}; it is exhausting for things like
\code{std::vector<std::pair<std::string, Element>>::const\_iterator}.
\cpp{}11 introduced \code{auto}: ``compiler, figure out the type from
the initialiser''.

\begin{lstlisting}
auto x = 42;                  // int
auto pi = 3.14;               // double
auto name = std::string{"hi"};// std::string
auto v = std::vector<int>{1, 2, 3};
for (auto& item : v) ++item;  // item is int&
\end{lstlisting}

Three rules of thumb:

\begin{itemize}[leftmargin=*]
  \item Plain \code{auto} drops references and \code{const}. If you
    write \code{auto x = some\_const\_ref();} you get a fresh copy.
  \item \code{auto\&} keeps references and \code{const}. Use this in
    range-\code{for} loops to avoid copying each element.
  \item \code{const auto\&} is the most common form for read-only
    iteration: \code{for (const auto\& cell : row) \{ ... \}}.
\end{itemize}

In \maya{}, \code{auto} is the default everywhere a type would be
verbose. Spell types out only when the spelling is short or when it
clarifies intent.

\section{Move semantics}
\label{sec:cpp:move}

We saw that assignment copies. For a \code{std::vector<int>} with a
million elements, copying is a million-element allocation plus a
million-element copy. If you do not need the original after, the
copy is pure waste.

Move semantics is the \cpp{}11 feature that lets you say ``transfer
the contents, do not copy''. A move is usually a pointer swap: the
new vector inherits the buffer; the old vector is left empty.

\begin{lstlisting}
std::vector<int> a = {1, 2, 3};
std::vector<int> b = std::move(a);   // O(1) --- b owns the buffer
// a is in a "valid but unspecified" state; do not read it.
\end{lstlisting}

Key points:

\begin{itemize}[leftmargin=*]
  \item \code{std::move(x)} does not actually move anything; it is a
    \emph{cast} that says ``treat \code{x} as a thing whose contents
    may be stolen''. The move actually happens in the constructor
    or assignment operator that receives the rvalue.
  \item After moving from \code{x}, you must not assume \code{x}
    holds anything specific. You can still destroy it (safely) and
    you can still assign a fresh value to it.
  \item Move is an \emph{optimisation}, not a semantic change. A
    correctly written program works whether \code{std::move} is
    present or not; \code{std::move} only changes how fast it runs.
\end{itemize}

The \textbf{rvalue reference} type \code{T\&\&} is how a function
opts in to receiving a move:

\begin{lstlisting}
void take(std::string s);              // by value: caller can move into it
void take_ref(const std::string& s);   // by const ref: read-only
void take_rref(std::string&& s);       // rvalue ref: must be moved in

std::string name = "alice";
take(name);              // copies name into the parameter
take(std::move(name));   // moves name into the parameter
// after this, name is moved-from
\end{lstlisting}

In \maya{} you see \code{std::move} most often in the runtime:
\code{Model} flows through \code{update(Model, Msg)} without being
copied, because \code{update} takes its \code{Model} by value and
the runtime moves the current model into the call.

\begin{warning}
Do not \code{std::move} a local on the return statement (\code{return
std::move(x);}). The compiler already moves return values
automatically; writing \code{std::move} explicitly disables a
separate optimisation called RVO (Return Value Optimisation) that is
usually \emph{better} than move.
\end{warning}

\section{Templates}
\label{sec:cpp:templates}

A \textbf{template} is a recipe the compiler runs at compile time.
You write the recipe once; the compiler bakes a fresh function or
class for each set of arguments you give it.

\subsection{Function templates}

\begin{lstlisting}
template <typename T>
T add(T a, T b) { return a + b; }

int    x = add(1, 2);          // T = int
double y = add(1.5, 2.5);      // T = double
std::string s = add<std::string>("foo", "bar");  // T = std::string
\end{lstlisting}

\code{typename T} introduces a type parameter named \code{T}. The
function body uses \code{T} as if it were a type. The compiler
\emph{instantiates} the template once per distinct \code{T}: there
are separate \code{add<int>}, \code{add<double>}, and
\code{add<std::string>} functions in the final binary.

The template only compiles correctly if \code{T} actually supports
\code{operator+}. Try \code{add(File, File)} and you will get a long
error. Concepts (Section~\ref{sec:cpp:concepts}) fix this.

\subsection{Class templates}

\begin{lstlisting}
template <typename T>
struct Box {
    T value;
    T& get() { return value; }
};

Box<int>    a{42};
Box<double> b{3.14};
Box<std::string> c{"hi"};
\end{lstlisting}

Each instantiation is a different type. \code{Box<int>} and
\code{Box<double>} share no relationship --- you cannot assign one
to the other.

\code{std::vector<T>}, \code{std::optional<T>},
\code{std::unique\_ptr<T>}, \code{std::variant<T, U, ...>} are all
class templates.

\subsection{Multiple template parameters}

\begin{lstlisting}
template <typename K, typename V>
struct Pair {
    K first;
    V second;
};

Pair<std::string, int> ages[] = {{"alice", 30}, {"bob", 25}};
\end{lstlisting}

\subsection{Non-type template parameters (NTTPs)}

The lesser-known but crucial flavour. A template parameter does not
have to be a type --- it can be a \emph{value}, computed at compile
time:

\begin{lstlisting}
template <int N>
struct Array {
    int data[N];
    int size() const { return N; }
};

Array<8>  a;   // data is int[8]
Array<16> b;   // data is int[16]
// a and b are different types
\end{lstlisting}

\code{std::array<T, N>} works this way. The size is part of the
type, so \code{std::array<int, 4>} and \code{std::array<int, 8>} are
distinct types and the size is known at compile time.

The types of NTTPs were limited for a long time --- integers,
pointers, enums. Since \cpp{}20, NTTPs can be \textbf{structural
class types}: any literal type with all-public members and no
user-defined comparison oddities. This is the trick that powers
\maya{}'s DSL.

\subsection{NTTPs of structural type --- the climax}

We want to write \code{t<"Hello">} and have the string become part of
the type. The straight path is to define a small ``string of fixed
length'' type and use it as an NTTP:

\begin{lstlisting}
template <std::size_t N>
struct Str {
    char data[N]{};

    consteval Str(const char (&s)[N]) {
        for (std::size_t i = 0; i < N; ++i) data[i] = s[i];
    }
};
\end{lstlisting}

Line by line:

\begin{itemize}[leftmargin=*]
  \item \code{template <std::size\_t N>}: a class template
    parameterised by an integer.
  \item \code{char data[N]\{\}}: a fixed-size array of \code{N}
    characters, zero-initialised.
  \item \code{consteval Str(const char (\&s)[N])}: a constructor that
    \emph{must} run at compile time (\code{consteval}) and takes a
    \emph{reference to a literal array of \code{N} chars}. The
    array-reference syntax \code{(\&s)[N]} captures the length of
    the literal --- the compiler picks \code{N} for you.
  \item The loop copies the string into \code{data}.
\end{itemize}

Now we can use \code{Str} as an NTTP:

\begin{lstlisting}
template <Str S>
struct Text {
    static constexpr std::string_view text() {
        return {S.data};
    }
};

Text<"Hello"> a;     // S = Str{"Hello\\0"}, baked into the type
Text<"World"> b;     // distinct type
\end{lstlisting}

\code{Text<"Hello">} and \code{Text<"World">} are different types.
The \emph{bytes} of the string are part of the type identity. This
is exactly how \maya{}'s \code{t<"...">} works, and it is the
foundation of the whole compile-time DSL we tour in
Chapter~\ref{ch:dsl}.

\begin{idea}
NTTPs let you encode \emph{values} into the type system. \maya{}
encodes text, colours, paddings, borders, and style flags this way.
Once a value is in a type, the compiler can check rules between
values at compile time.
\end{idea}

\subsection{Templates and the \code{<>} syntax}

When you see \code{Foo<X, Y>} in a \cpp{} expression, the angle
brackets are template-argument syntax, not less-than / greater-than.
The ambiguity is one of the language's papercuts; the compiler
usually figures it out from context, and where it cannot, you write
\code{template} explicitly: \code{x.template foo<int>()}.

You will see \code{t<"Hi">}, \code{Fg<255, 100, 100>}, \code{pad<1>},
\code{border\_<Round>} throughout \maya{}. They are all template
instantiations.

\section{\code{constexpr}, \code{consteval}, \code{constinit}}
\label{sec:cpp:constexpr}

A \textbf{constant expression} is one the compiler can evaluate at
translation time, before the program runs. \cpp{} has three
keywords governing this:

\begin{itemize}[leftmargin=*]
  \item \code{constexpr} on a function or variable: \emph{may} be
    constant-evaluated; otherwise behaves like a normal function or
    variable.
  \item \code{consteval} on a function: \emph{must} be
    constant-evaluated. Calling it at runtime is a compile error.
  \item \code{constinit} on a variable: must be initialised by a
    constant expression, but the variable itself is otherwise
    mutable. Used to forbid the ``static initialisation order
    fiasco''.
\end{itemize}

Example of \code{constexpr}:

\begin{lstlisting}
constexpr int square(int n) { return n * n; }

constexpr int sixteen = square(4);   // evaluated at compile time
int runtime_n = 7;
int forty_nine = square(runtime_n);  // also legal, evaluated at runtime
\end{lstlisting}

Example of \code{consteval}:

\begin{lstlisting}
consteval int square_only_ct(int n) { return n * n; }

constexpr int a = square_only_ct(4);  // ok
int runtime_n = 7;
// int b = square_only_ct(runtime_n); // error: must be constant-evaluated
\end{lstlisting}

\maya{}'s DSL is heavy on \code{constexpr} and \code{consteval}. The
entire chain \code{t<"Hi"> | Bold | Fg<255, 80, 80>} produces a
\code{constexpr} value --- every \code{operator|} along the way is
\code{constexpr}. Conversion to the runtime \code{Element} (which
may allocate) happens only when you call \code{.build()}.

\section{Concepts and \code{requires}}
\label{sec:cpp:concepts}

A template constrains its arguments implicitly: writing \code{a + b}
in a template body means \code{T} must support \code{operator+}, but
the constraint is invisible until you try to instantiate the
template with a wrong type and get a wall of compiler errors.

\cpp{}20 \textbf{concepts} make those constraints explicit and
named:

\begin{lstlisting}
template <typename T>
concept Addable = requires(T a, T b) {
    { a + b } -> std::convertible_to<T>;
};

template <Addable T>
T add(T a, T b) { return a + b; }

add(1, 2);                // ok, int is Addable
// add(File{...}, File{...});  // clean error: File does not satisfy Addable
\end{lstlisting}

The \code{requires} expression lists statements that must be valid
for \code{T}. \code{\{ a + b \} -> std::convertible\_to<T>} means
``\code{a + b} must compile, and its result must convert to
\code{T}''. If a caller passes a \code{T} that does not satisfy
\code{Addable}, the compiler points at the constraint, not at
something seventeen layers deep inside the template body.

\subsection{Maya's concepts}

Maya defines several concepts in \file{core/concepts.hpp}. Three of
the important ones:

\begin{lstlisting}
// (paraphrased from maya source)

template <typename T>
concept Node = requires(const T& t) {
    { t.build() } -> std::convertible_to<Element>;
};

template <typename P>
concept Program = requires(typename P::Model m, typename P::Msg msg) {
    typename P::Model;
    typename P::Msg;
    { P::init() };
    { P::update(m, msg) };
    { P::view(m) }       -> std::convertible_to<Element>;
    { P::subscribe(m) };
};
\end{lstlisting}

\code{Node} says ``has a \code{.build()} that returns something
convertible to \code{Element}''. Every DSL piece satisfies it.
\code{Program} says ``has the four Elm functions with the right
signatures''. \code{run<P>()} requires \code{Program<P>}; if your
program is missing a function, the compiler tells you which one.

\begin{idea}
Concepts are documentation that the compiler checks. They turn
``this template requires the type to do X'' from a comment into a
compile-time predicate.
\end{idea}

\section{\code{std::variant} and \code{std::visit}}
\label{sec:cpp:variant}

A \code{std::variant<A, B, C>} is a value that holds \emph{exactly
one} of \code{A}, \code{B}, or \code{C}. It is the \cpp{} closed sum
type --- the same construct functional languages call an algebraic
data type.

\begin{lstlisting}
#include <variant>
std::variant<int, std::string, double> v;
v = 42;          // holds an int
v = "hello";     // now holds a string
v = 3.14;        // now holds a double
\end{lstlisting}

You inspect a variant with \code{std::visit}, which dispatches to
the overload matching the current alternative. The idiomatic helper
is \code{overload}, a small struct that combines multiple lambdas
into a single callable:

\begin{lstlisting}
template <typename... Fs> struct overload : Fs... { using Fs::operator()...; };
template <typename... Fs> overload(Fs...) -> overload<Fs...>;

std::visit(overload{
    [](int n)               { std::println("int  : {}", n); },
    [](const std::string& s){ std::println("str  : {}", s); },
    [](double d)            { std::println("real : {}", d); },
}, v);
\end{lstlisting}

The \code{overload} template uses two features:

\begin{itemize}[leftmargin=*]
  \item \code{Fs...} is a \textbf{parameter pack} --- a variadic
    template, accepting any number of types.
  \item Inheriting from each one and using \code{using
    Fs::operator()...} brings every lambda's \code{operator()} into
    the combined class. \code{overload\{f, g, h\}} is then a single
    object with three call-operator overloads.
\end{itemize}

This is enough \cpp{} sorcery to write off as a one-liner you copy.
It is in \file{maya/core/overload.hpp}.

\subsection{Maya's variants}

Maya uses \code{std::variant} relentlessly:

\begin{itemize}[leftmargin=*]
  \item \code{Element} is
    \code{std::variant<BoxElement, TextElement, ElementList>}.
  \item \code{Event} is
    \code{std::variant<Key, Mouse, Resize, Paste, ...>}.
  \item \code{Cmd<Msg>} is a variant over every effect the runtime
    knows: \code{None}, \code{Quit}, \code{Batch}, \code{After},
    \code{SetTitle}, \code{WriteClipboard}, \code{Task}, ...
  \item \code{Sub<Msg>} is a variant over event sources.
\end{itemize}

Whenever you see a \code{std::visit} in \maya{}, it is the dispatch
hub for one of these types.

\begin{warning}
If you add a new alternative to a variant and forget to handle it in
a visitor, the visitor stops compiling. This is a feature: it makes
refactors safe. Embrace it; do not paper over it with a generic
fallback.
\end{warning}

\section{\code{std::expected<T, E>}}
\label{sec:cpp:expected}

\cpp{}23's \code{std::expected<T, E>} is a sum type for fallible
operations: either a success of type \code{T} or an error of type
\code{E}. It is \maya{}'s preferred alternative to throwing
exceptions for ``failure is expected, just rare'' cases.

\begin{lstlisting}
#include <expected>
#include <charconv>

std::expected<int, std::string> parse_int(std::string_view s) {
    int n;
    auto [end, ec] = std::from_chars(s.data(), s.data() + s.size(), n);
    if (ec != std::errc{}) return std::unexpected("not an int");
    return n;
}

if (auto r = parse_int("42")) {
    std::println("value: {}", *r);
} else {
    std::println("error: {}", r.error());
}
\end{lstlisting}

The API at a glance:

\begin{itemize}[leftmargin=*]
  \item \code{(bool)r} --- true if success.
  \item \code{*r} or \code{r.value()} --- the success value.
  \item \code{r.error()} --- the error value.
  \item \code{std::unexpected(e)} --- factory for the error case.
\end{itemize}

In \maya{} the error type is usually a small struct:

\begin{lstlisting}
enum class HttpErrorKind { Tls, Timeout, Refused, Status, Cancel };
struct HttpError {
    HttpErrorKind kind;
    int           http_status;
    std::string   detail;
};
using HttpResult = std::expected<Response, HttpError>;
\end{lstlisting}

The caller switches on \code{kind} without parsing strings --- a
proper closed sum of failure modes. This is what we call a
``refined'' error type.

\section{Lambdas and capture}
\label{sec:cpp:lambdas}

A \textbf{lambda} is an anonymous function object. Under the hood,
the compiler synthesises a nameless class with an \code{operator()}
that runs the lambda body.

\begin{lstlisting}
auto add = [](int a, int b) { return a + b; };
int n = add(3, 4);    // 7
\end{lstlisting}

The square brackets at the front are the \textbf{capture list}: they
list names from the surrounding scope the lambda should remember.

\begin{lstlisting}
int x = 10;

auto by_value = [x](int n) { return x + n; };   // copy of x
auto by_ref   = [&x](int n){ return x + n; };   // reference to x

x = 99;
by_value(0);   // 10  --- the lambda has the old copy
by_ref(0);     // 99  --- the lambda sees the new x
\end{lstlisting}

Shorthand captures:

\begin{lstlisting}
[=]        // capture every used name by value
[&]        // capture every used name by reference
[=, &out]  // mostly by value, but out by reference
[&, n]     // mostly by reference, but n by value
\end{lstlisting}

Default captures (\code{[=]} and \code{[\&]}) are convenient but
dangerous in long-lived lambdas; you can easily capture a reference
to a local that has gone out of scope. Prefer explicit capture lists.

\subsection{Init-captures and moving into a lambda}

\cpp{}14 added \textbf{init-captures}: \code{[name = expression]}
introduces a new name local to the lambda with any initialiser. This
is how you capture a move-only type into a lambda:

\begin{lstlisting}
auto buf = std::make_unique<Buffer>();
auto cb  = [buf = std::move(buf)] {
    buf->flush();
};
// outer buf is moved-from; only the lambda owns the buffer.
\end{lstlisting}

\subsection{Storing lambdas: \code{std::function} vs.\ \code{std::move\_only\_function}}

A lambda's type is unique to it; you cannot write the type by hand.
To store a lambda in a class member or container, you need a
type-erased wrapper:

\begin{itemize}[leftmargin=*]
  \item \code{std::function<R(Args...)>} (\cpp{}11) holds any
    callable matching the signature, but the callable must be
    \emph{copyable}. That excludes lambdas capturing move-only state
    (\code{std::unique\_ptr}, sockets, file handles).
  \item \code{std::move\_only\_function<R(Args...)>} (\cpp{}23)
    requires only \emph{movable}. \maya{}'s signal callbacks and
    \code{Cmd<Msg>::Task} use this; effect descriptions routinely
    carry move-only state.
\end{itemize}

\section{Standard-library bits you will see}
\label{sec:cpp:stdlib}

The standard library is large. A non-exhaustive list of names that
appear in \maya{} and what they are:

\begin{itemize}[leftmargin=*]
  \item \code{std::string\_view} --- a non-owning view of a
    contiguous \code{char} sequence. Two pointers under the hood.
    Cheap to pass; does not own the bytes; the underlying data must
    outlive the view.
  \item \code{std::span<T>} --- the same idea for any contiguous
    range of \code{T}. \code{std::span<const Cell>} is how the
    renderer receives a row.
  \item \code{std::array<T, N>} --- a fixed-size array; size is a
    template argument. Same memory layout as \code{T[N]} but with
    proper value semantics and methods.
  \item \code{std::vector<T>} --- the dynamic-array workhorse.
    Heap-allocated buffer that grows.
  \item \code{std::optional<T>} --- either a \code{T} or nothing. Use
    when ``absent'' is a real state but you do not need a richer
    error.
  \item \code{std::format} / \code{std::print} / \code{std::println}
    --- \cpp{}20/23 formatted output, like Python's f-strings.
  \item \code{std::source\_location} --- a value carrying the file,
    line, and function name where it was constructed. \maya{} uses
    it for stable cache identifiers.
  \item \code{std::chrono} --- duration and time-point types.
    \code{std::chrono::milliseconds(50)} reads cleanly.
  \item \code{std::atomic<T>} --- lock-free operations on a value
    shared across threads.
  \item \code{std::jthread} --- a thread that joins on destruction
    (no leaked threads on early return). \cpp{}20.
  \item \code{std::ranges} / \code{std::views} --- lazy, composable
    algorithms over sequences (\code{xs | views::transform(...) |
    views::filter(...)}).
\end{itemize}

Learn the views as you meet them; the ranges library is large and you
do not need it all at once.

\section{Reading real \maya{} code}
\label{sec:cpp:reading}

With the toolkit in hand, let us read a fragment of actual \maya{}.
This is from the style header; it declares the tag types behind
\code{Bold}, \code{Fg<R, G, B>}, and friends.

\begin{lstlisting}
template <CTStyle V>
struct StyTag {
    constexpr operator Style() const noexcept { return V.runtime(); }
};

inline constexpr StyTag<CTStyle::bold()>   Bold{};
inline constexpr StyTag<CTStyle::italic()> Italic{};

template <std::uint8_t R, std::uint8_t G, std::uint8_t B>
inline constexpr StyTag<CTStyle::fg(R, G, B)> Fg{};
\end{lstlisting}

Line by line:

\begin{itemize}[leftmargin=*]
  \item \code{template <CTStyle V>}: \code{StyTag} is a class
    template parameterised by an NTTP \code{V} of type
    \code{CTStyle}. (\code{CTStyle} is a structural class type
    --- another struct defined elsewhere --- carrying a compile-time
    description of a style: which attributes are on, the colours,
    and so on.)
  \item \code{struct StyTag \{ ... \}}: the type has no data
    members. It is a \emph{tag} type --- empty at runtime, distinct
    at compile time. Two \code{StyTag<X>}'s with different \code{X}
    are different types.
  \item \code{constexpr operator Style() const noexcept}: an
    implicit conversion from \code{StyTag<V>} to the runtime
    \code{Style} type. \code{constexpr}, so it can run at compile
    time; \code{noexcept} promises it does not throw. \code{const}
    on the trailing position says the conversion does not modify
    the object.
  \item \code{return V.runtime();}: ask the compile-time style for
    its runtime form. \code{V} is the NTTP, available inside the
    class.
  \item \code{inline constexpr StyTag<CTStyle::bold()> Bold\{\};}: a
    \code{constexpr} variable of type
    \code{StyTag<CTStyle::bold()>}, initialised with default braces.
    \code{inline} on a variable in a header lets it appear in
    multiple translation units without violating ODR. \code{Bold} is
    the value you write in user code.
  \item \code{Fg<R, G, B>}: a \emph{variable template} ---  a
    template that yields a variable, not a function or class. The
    expression \code{Fg<255, 80, 80>} instantiates the template with
    those three integers, producing a unique \code{StyTag<...>}
    value.
\end{itemize}

If you can read that fragment now and feel the pieces in place ---
NTTP, class template, \code{constexpr}, \code{operator T()},
variable template, the implicit conversion that lets \code{Bold}
appear in places that want a \code{Style} --- you can read the rest
of \maya{}'s headers. From here on, the rest of the book is mostly
about what each file \emph{does}, not how its syntax decodes.

\begin{recap}
The \cpp{} \maya{} uses, in one paragraph: templates with type and
non-type parameters; \code{constexpr}/\code{consteval} for compile-
time evaluation; concepts and \code{requires} for named constraints;
\code{std::variant} + \code{std::visit} + \code{overload} for closed
sum types; \code{std::expected} for structured errors; lambdas with
explicit capture and \code{std::move\_only\_function} for storage;
move semantics to avoid copies; RAII to make resources self-managing.
Together they support one discipline: encode invariants in the type
system, make impossible states unrepresentable, defer side effects to
the edge.
\end{recap}

\section*{Workshop}

The exercises here are calibrated for someone who has never used the
feature. If a step takes you fifteen minutes and three Google
searches, that is the right speed.

\begin{enumerate}[leftmargin=*]
  \item \textbf{RAII guard.} Write a class \code{Timer} whose
    constructor records \code{std::chrono::steady\_clock::now()} and
    whose destructor prints the elapsed time. Use it at the top of
    a function; verify the message appears on return.
  \item \textbf{Templates and NTTPs.} Implement a \code{FixedArray<T,
    N>} class template with a member \code{T data[N]}, a
    \code{size()} method returning \code{N}, and \code{operator[]}.
    Instantiate \code{FixedArray<int, 4>} and \code{FixedArray<int,
    8>}; convince yourself they are different types by trying to
    assign one to the other.
  \item \textbf{NTTP of structural type.} Type out the \code{Str<N>}
    template from this chapter. Define a class template \code{Tag<Str
    S>} parameterised by it. Instantiate \code{Tag<"hello">} and
    \code{Tag<"world">} and call a method that returns the string.
  \item \textbf{Concept.} Write a concept \code{Drawable} requiring
    a \code{void draw() const} method. Use it on a template function
    \code{render(const T\& t)}. Make two structs satisfying it, one
    not, and observe the error.
  \item \textbf{Variant.} Define
    \code{using Shape = std::variant<Circle, Square, Triangle>}
    with three trivial structs (each with one field). Write a
    \code{double area(const Shape\&)} using \code{std::visit} and the
    \code{overload} helper.
  \item \textbf{Expected.} Write
    \code{std::expected<int, std::string> safe\_div(int a, int b)}.
    Call it with \code{(10, 2)} and \code{(10, 0)}; print each
    result.
  \item \textbf{Lambda with move capture.} Build a
    \code{std::vector<int>}, capture it into a lambda by move, and
    return the lambda from a function. Call the returned lambda;
    confirm the vector outlives the function it came from.
\end{enumerate}
